% Winkel und Dreiecke – bank/winkel-dreiecke/ (zusammenbau v0.4, Heft)
\einheitenkopf[t7]{Winkel und Dreiecke}
\setcounter{aufgabe}{47}

% Anlauf (ohne Prüfkennung)
\begin{aufgabe}{Winkel messen – Anlauf}
\begin{teile}
\teil Welche Winkelart hat α? Nicht messen. \\ \kreuz{spitz} \\ \kreuz{recht} \\ \kreuz{stumpf} \\ \kreuz{gestreckt} \\ \kreuz{überstumpf}

\winkel[4]{27}{\alpha}
\teil Miss den Winkel α mit dem Geodreieck.

\winkel[5]{37}{\alpha}
α = \leerfeld °
\teil Miss den stumpfen Winkel α mit dem Geodreieck.

\winkel[5]{118}{\alpha}
α = \leerfeld °
\end{teile}
\end{aufgabe}

\begin{aufgabe}{Winkel messen – Prüfungsaufgaben}
\begin{teile}
\teil Vom Punkt A aus sieht Jana die Spitze T eines Turms unter einem Winkel von $41^\circ$ zur waagerechten Grundlinie, die Spitze F der Fahne auf dem Turm unter $44^\circ$. Wie groß ist der Winkel γ zwischen den Sichtlinien AT und AF? \hfill (P10 2018 OS) \\ γ = \leerfeld °
\teil Von der Kaimauer aus bildet die Sichtlinie zur Lampe eines Leuchtturms mit der Wasserlinie einen Winkel von $22^\circ$, die Sichtlinie zur Spitze eines Krans dahinter einen Winkel von $29^\circ$. Wie groß ist der Winkel γ zwischen den beiden Sichtlinien? \hfill (P10 2018 OS) \\ γ = \leerfeld °
\teil Ein Radweg führt von P über Q nach R. Die Strecke PQ ist 730 m lang, die Strecke QR 684 m. Wie lang ist der Radweg von P nach R? \hfill (P10 2016 OS) \\ \leerfeld[m]
\teil Pauls Schulweg führt von der Haustür 615 m bis zur Ampel und von dort 470 m bis zur Schule. Wie lang ist sein Schulweg? \hfill (P10 2016 OS) \\ \leerfeld[m]
\end{teile}
\end{aufgabe}

% Anlauf (ohne Prüfkennung)
\begin{aufgabe}{Winkel an Geraden – Anlauf}
\begin{teile}
\teil Wie liegen α und γ an der Kreuzung zueinander? \\ \kreuz{Scheitelwinkel} \\ \kreuz{Nebenwinkel} \\ \kreuz{Stufenwinkel} \\ \kreuz{Wechselwinkel}

\geradenkreuzung{55}{\alpha}{\beta}{\gamma}{\delta}
\teil An einer Geradenkreuzung ist $\alpha = 38^\circ$. Wie groß ist der Scheitelwinkel von α?

\geradenkreuzung{38}{\alpha}{}{}{}
\leerfeld °
\teil An einer Geradenkreuzung ist $\alpha = 44^\circ$. Wie groß ist ein Nebenwinkel von α?

\geradenkreuzung{44}{\alpha}{}{}{}
\leerfeld °
\end{teile}
\end{aufgabe}

\begin{aufgabe}{Winkel an Geraden – Prüfungsaufgaben}
\begin{teile}
\teil Die Geraden g und h schneiden sich; zwischen ihnen liegt $\alpha = 64^\circ$. Eine dritte Gerade durch den Schnittpunkt teilt den Scheitelwinkel von α in zwei Teile, der untere misst $23^\circ$. Wie groß ist der obere Teil β? \hfill (P10 2014 OS) \\ β = \leerfeld °
\teil $g \parallel h$, g liegt oben, h unten, die Gerade k schneidet beide. An g liegt unterhalb von g links von k ein Winkel von $49^\circ$. Wie groß ist der Winkel an h oberhalb von h rechts von k? \hfill (P10 2020 OS) \\

\parallelenpaar{49}{}{}{}{}
\leerfeld °
\teil Im Dreieck ABC ist der Winkel bei C $58^\circ$ groß. Die Seiten AB und CB werden über B hinaus verlängert; zwischen den beiden Verlängerungen liegt ein Winkel von $43^\circ$. Wie groß ist der Innenwinkel β bei B? \hfill (P10 2019 OS) \\ β = \leerfeld °
\teil Im Trapez ABCD mit $AB \parallel CD$ ist $\beta = 58^\circ$, $\gamma = 122^\circ$ und $\delta = 107^\circ$. Wie groß ist α? \hfill (P10 2021 OS) \\

\viereck[winkel]{(0,0)}{(6,0)}{(4.44,2.5)}{(0.76,2.5)}
α = \leerfeld °
\teil Ein Deich hat im Querschnitt die Form eines symmetrischen Trapezes, die lange Seite liegt unten. Die Böschung steigt an beiden Seiten unter $52^\circ$ an. Wie groß ist der Winkel γ oben rechts an der Deichkrone? \hfill (P10 2023 OS) \\

\trapez{5}{3}{1.28}
γ = \leerfeld °
\end{teile}
\end{aufgabe}

\begin{aufgabe}{Winkel an Geraden – Prüfungsaufgaben – weiter}
\begin{teile}
\teil Die parallelen Geraden h (oben) und g (unten) werden von der Geraden k geschnitten. An h liegt unterhalb von h links von k ein Winkel von $44^\circ$; α liegt an g oberhalb von g links von k. Wie groß ist α? \hfill (P10 2015 OS) \\

\parallelenpaar{44}{}{}{}{}
α = \leerfeld °
\teil Eine gerade Straße k überquert zwei parallele Gleise, in der Skizze liegt Gleis h oben, Gleis g unten. Unterhalb von h links von k bilden Straße und Gleis einen Winkel von $51^\circ$. Wie groß ist der Winkel α an g oberhalb von g links von k? \hfill (P10 2015 OS) \\

\parallelenpaar{51}{}{}{}{}
α = \leerfeld °
\teil Ein Parallelogramm hat links unten einen Winkel von $79^\circ$. Wie groß ist der benachbarte Winkel β rechts unten? \hfill (P10 2026 FOR) \\

\parallelogramm[winkel={79^\circ,\beta,,}]{4}{2.5}{79}
β = \leerfeld °
\teil Ein Scherengitter ist aus Parallelogrammen aufgebaut. In einem davon misst der Winkel links unten $83^\circ$. Wie groß ist der benachbarte Winkel β rechts unten? \hfill (P10 2026 FOR) \\

\parallelogramm[winkel={83^\circ,\beta,,}]{4}{2.5}{83}
β = \leerfeld °
\end{teile}
\end{aufgabe}

% Anlauf (ohne Prüfkennung)
\begin{aufgabe}{Winkelsumme – Anlauf}
\begin{teile}
\teil Markiere gleich lange Seiten oder den rechten Winkel und kreuze die Dreiecksart an. \\ \kreuz{spitzwinklig} \\ \kreuz{rechtwinklig} \\ \kreuz{stumpfwinklig}

\dreieckrw{4}{3}{a}{b}{c}
\teil Im Dreieck ABC ist $\alpha = 52^\circ$ und $\beta = 71^\circ$. Wie groß ist γ? γ = \leerfeld °
\teil Im Dreieck ABC ist $\alpha = 38{,}6^\circ$ und $\beta = 94{,}7^\circ$. Wie groß ist γ? \leerfeld °
\end{teile}
\end{aufgabe}

\begin{aufgabe}{Winkelsumme – Prüfungsaufgaben}
\begin{teile}
\teil Ein dreieckiges Grundstück hat an zwei Ecken die Winkel $51{,}4^\circ$ und $76{,}9^\circ$. Wie groß ist der Winkel an der dritten Ecke? \hfill (P10 2017 OS) \\ \leerfeld °
\teil Im Giebeldreieck ABC eines Hauses steht die Stütze CD senkrecht auf dem Balken AB. An der Ecke A ist $\alpha = 27{,}6^\circ$. Wie groß ist der Winkel $\gamma_1$ zwischen AC und der Stütze CD? \hfill (P10 2022 OS) \\ γ₁ = \leerfeld °
\teil Die Punkte B, C und D liegen senkrecht übereinander, C zwischen B und D. Im Dreieck ABC ist der Winkel bei A $22^\circ$ groß und der Winkel bei B ein rechter. Der Winkel zwischen AC und AD beträgt $9^\circ$. Zeige, dass der Winkel δ bei D im Dreieck ACD $59^\circ$ groß ist. \hfill (P10 2018 OS)
\teil Im Drachenviereck ABCD ist BD die Symmetrieachse. Im Dreieck ABD ist $\angle DAB = 114^\circ$ und $\angle ABD = 38^\circ$. Wie groß ist $\delta = \angle ADB$? \hfill (P10 2015 OS) \\ δ = \leerfeld °
\teil Im Dachdreieck ABC steht die Stütze BF senkrecht auf dem Balken AC, F liegt auf AC. Die Dachschräge AB ist gegen den Balken um einen halben rechten Winkel geneigt. Begründe, dass die Stütze BF genauso lang ist wie das Balkenstück AF. \hfill (P10 2023 OS)
\teil Ergänze richtig: „In jedem Drachenviereck …“ \hfill (P10 2026 FOR) \\ \kreuz{sind alle Winkel gleich groß} \\ \kreuz{gibt es zwei Paare gleich langer Nachbarseiten} \\ \kreuz{sind alle Seiten gleich lang}
\teil Ergänze richtig: „In jeder Raute …“ \hfill (P10 2016 OS) \\ \kreuz{sind benachbarte Winkel gleich groß} \\ \kreuz{sind gegenüberliegende Winkel gleich groß} \\ \kreuz{sind alle Winkel gleich groß}
\end{teile}
\end{aufgabe}

\begin{aufgabe}{Winkelsumme – Prüfungsaufgaben – weiter}
\begin{teile}
\teil Im Dreieck ABC ist $\alpha = \beta = 66^\circ$ und AC = 7 cm. Gib die Länge von BC und die Größe von γ an. \hfill (P10 2023 OS) \\ BC = \leerfeld[cm] , γ = \leerfeld °
\teil Die Giebelwand eines Zeltes ist ein Dreieck ABC mit den Winkeln $\alpha = \beta = 73^\circ$ an den Bodenecken; die Kante AC ist 2,4 m lang. Wie lang ist die Kante BC, und wie groß ist der Winkel γ an der Spitze? \hfill (P10 2023 OS) \\ BC = \leerfeld[m] , γ = \leerfeld °
\steil Im Drachenviereck ABCD schneiden sich die Diagonalen AC und BD im Punkt F senkrecht. F ist die Mitte von AC, AC = 36 cm und DF = 18 cm. Begründe, dass der Winkel des Drachens bei D ein rechter ist. \hfill (P10 2025 OS)

\drachen[diagonalen]{5}{3.6}{0.36}
\steil Bei einem Flugdrachen ABCD wird die 44 cm lange Querstange AC von der Längsstange BD in ihrer Mitte F senkrecht gekreuzt. Die Spitze D ist 22 cm von F entfernt. Begründe, dass die oberen Kanten AD und CD bei D einen rechten Winkel bilden. \hfill (P10 2025 OS)
\end{teile}
\end{aufgabe}

% Anlauf (ohne Prüfkennung)
\begin{aufgabe}{Konstruieren – Anlauf}
\begin{teile}
\teil Gegeben sind a = 6 cm, b = 4 cm, $\gamma = 52^\circ$. Markiere die Stücke in einer Planfigur. Welcher Kongruenzsatz passt? \\ \kreuz{SSS} \\ \kreuz{SWS} \\ \kreuz{WSW} \\ \kreuz{SsW}
\teil Konstruiere das Dreieck ABC mit a = 6 cm, b = 5 cm und c = 7 cm. Beginne mit der Seite c auf dem Strahl.

\winkelstrahl[9][A]
\teil Konstruiere das Dreieck ABC mit b = 5 cm, c = 6 cm und $\alpha = 50^\circ$. Beginne mit der Seite c auf dem Strahl.

\winkelstrahl[9][A]
\end{teile}
\end{aufgabe}

\begin{aufgabe}{Konstruieren – Prüfungsaufgaben}
\begin{teile}
\teil Ein dreieckiges Beet ABC hat die Seite a = 12 m. Was gilt für die Summe der beiden anderen Seiten b und c? \hfill (P10 2020 OS) \\ \kreuz{$b + c > 12$ m} \\ \kreuz{$b + c = 12$ m} \\ \kreuz{$b + c < 12$ m}
\teil Drei Holzstäbe sollen zu einem Dreieck gelegt werden; der längste Stab a ist 15 cm lang. Was muss für die beiden anderen Stäbe b und c gelten? \hfill (P10 2020 OS) \\ \kreuz{$b + c > 15$ cm} \\ \kreuz{$b + c = 15$ cm} \\ \kreuz{$b + c < 15$ cm}
\end{teile}
\end{aufgabe}
